\honest{Superexponential Conceptspace, and Simple Words}{Eliezer Yudkowsky}{2008}

Thingspace holds everything that could exist. I admit that, the way I defined it, Thingspace ``may be too poorly defined to have anything you could call a size.'' Still, I say, the space of concepts is far larger. \nb{That holds whatever size Thingspace has: a concept is a set of things, and there are always more sets of things than things.}

In machine learning, a concept is a rule that includes or excludes examples. My example, adapted from Tom Mitchell's textbook \textsc{Machine Learning}, is days for playing tennis, described by Sky (three values) and by AirTemp, Humidity and Wind (two values each). Given two good days and one bad one, a learner might guess \{?, Warm, High, ?\}: warm and humid, with any sky and any wind. The classic algorithm keeps two sets, the most general and the most specific concepts that fit the data, and moves them toward each other as examples come in. Here the most general are \{?, Warm, ?, ?\} and \{Sunny, ?, ?, ?\}, and the most specific is \{Sunny, Warm, High, ?\}. \nb{All of this checks. Of the 109 concepts in this format, six fit the three examples, and the two sets are as stated.}

This format cannot express ``Play tennis when the sky is sunny or the air is warm.'' So why not let the learner consider every concept? There are 24 possible days and 109 concepts in the format, but every set of days is a concept, and there are $2^{24} = 16{,}777{,}216$ sets of days. A learner that allows all of them, I say, ``can't learn concepts until they've seen every possible example.'' \nb{That is correct: for any day not yet seen, exactly half of the concepts that fit the data accept it. Mitchell's lecture slides for that chapter make the same point.}

A restricted space, on the other hand, can be narrowed with a number of examples that grows only with the logarithm of its size. Add a Water attribute, and there are 48 days and 325 concepts. If each example is accepted by about half of the concepts still in play, nine examples will do; with forty yes-or-no attributes, 64. \nb{In this format no first example comes near half: any one day is accepted by 32 of the 325 concepts. When each example was chosen to split the remaining concepts as evenly as possible, a simulation needed about 15 on average. That the number grows with the logarithm is a standard result.} The space of all concepts over forty attributes, though, has $2^{2^{40}}$ members, and to learn in it you would have to see every one of the trillion or so possible examples. So real minds think only about highly regular concepts, and learning ``is nearly all inductive bias.'' \nb{In the four-attribute example, the format rules out 16,777,107 of the 16,777,216 concepts before any data, and the three examples rule out 103 more.}

What has this to do with words? ``It's the whole reason that words have intensions as well as extensions.'' Yesterday I said that the way to carve reality at its joints is to draw boundaries around concentrations of unusually high probability density. I left out one word on purpose: simple. Without it you would gerrymander Thingspace, drawing a boundary that only lists what you have seen. Socrates is not shaped exactly like any other human, so calling him human needs a simple boundary, not a list of five-megabyte shape specifications.

I admit that the rule is ``a bit chicken-and-eggish'': you cannot judge density before you have done some carving, and the probabilities come from the boundaries. If you already had the probabilities, I ask, ``why would you bother drawing boundaries?'' \nb{Yudkowsky gave a reason a week later, in ``Conditional Independence, and Naive Bayes'': pretending that one central class variable screens off the others ``can simplify the living daylights out of your calculations.''} I go on to my next point without resolving the circle.

Given how many concepts there are, singling out one for attention ``is an act of no small audacity.'' Defining ``wiggin'' as ``a black-haired green-eyed person'' with no reason to attend to it is like a detective who, with no evidence at all, asks whether we have considered John Q. Wiffleheim of 1234 Norkle Rd. \nb{``Wiggin'' is a conjunction of two attribute values, one of the simple concepts of the restricted format, so the superexponential count is not what makes it audacious. The detective's point holds in the smaller space as well.}
